\documentclass[11pt]{article}

\usepackage[T1]{fontenc}
\usepackage[utf8]{inputenc}
\usepackage{lmodern}
\usepackage{microtype}
\usepackage[margin=1in,headheight=14pt]{geometry}
\usepackage{amsmath,amssymb,amsthm,mathtools}
\usepackage{mathrsfs}
\usepackage{booktabs,longtable,tabularx,array}
\usepackage{enumitem}
\usepackage{xcolor}
\usepackage[most]{tcolorbox}
\usepackage{fancyhdr}
\usepackage{xurl}
\usepackage{hyperref}
\usepackage[nameinlink,capitalise,noabbrev]{cleveref}
\usepackage{listings}
\usepackage{setspace}

\definecolor{navy}{RGB}{20,49,83}
\definecolor{bluegray}{RGB}{67,87,111}
\definecolor{newgreen}{RGB}{19,105,71}
\definecolor{openorange}{RGB}{163,83,20}
\definecolor{softblue}{RGB}{240,246,252}
\definecolor{softgreen}{RGB}{240,249,245}
\definecolor{softorange}{RGB}{253,247,239}
\definecolor{softgray}{RGB}{246,247,249}
\definecolor{deepred}{RGB}{143,36,36}

\hypersetup{
  colorlinks=true,
  linkcolor=navy,
  citecolor=newgreen,
  urlcolor=bluegray,
  pdftitle={Alaoglu--Erdos: Dyadic Cyclotomic Defects and Fresh Prime Supports},
  pdfauthor={Research continuation prepared by GPT-5.6 Pro},
  pdfsubject={Rigorous partial progress on the Alaoglu--Erdos integer-exponent conjecture}
}

\pagestyle{fancy}
\fancyhf{}
\fancyhead[L]{\small Alaoglu--Erd\H{o}s continuation report}
\fancyhead[R]{\small 22 August 2026}
\fancyfoot[C]{\thepage}

\setlength{\parindent}{0pt}
\setlength{\parskip}{0.55em}
\setlist[itemize]{leftmargin=1.6em,itemsep=0.25em,topsep=0.25em}
\setlist[enumerate]{leftmargin=1.8em,itemsep=0.25em,topsep=0.25em}
\renewcommand{\arraystretch}{1.15}
\allowdisplaybreaks
\setlength{\emergencystretch}{2em}

\newtheorem{theorem}{Theorem}[section]
\newtheorem{proposition}[theorem]{Proposition}
\newtheorem{corollary}[theorem]{Corollary}
\newtheorem{lemma}[theorem]{Lemma}
\newtheorem{conjecture}[theorem]{Target conjecture}
\theoremstyle{definition}
\newtheorem{definition}[theorem]{Definition}
\newtheorem{remark}[theorem]{Remark}
\newtheorem{audit}[theorem]{Failure audit}

\newtcolorbox{aeexecutivebox}[1][]{
  colback=softblue,colframe=navy,boxrule=0.8pt,arc=2pt,
  title=Executive assessment,fonttitle=\bfseries,#1}
\newtcolorbox{aenovelbox}[1][]{
  colback=softgreen,colframe=newgreen,boxrule=0.8pt,arc=2pt,
  title=New rigorous result in this continuation,fonttitle=\bfseries,#1}
\newtcolorbox{aeopenbox}[1][]{
  colback=softorange,colframe=openorange,boxrule=0.8pt,arc=2pt,
  title=Open closing step,fonttitle=\bfseries,#1}
\newtcolorbox{aestatusbox}[1][]{
  colback=softgray,colframe=bluegray,boxrule=0.6pt,arc=2pt,
  title=Status and trust boundary,fonttitle=\bfseries,#1}
\newtcolorbox{aewarningbox}[1][]{
  colback=red!3,colframe=deepred,boxrule=0.8pt,arc=2pt,
  title=Do not overclaim,fonttitle=\bfseries,#1}

\newcommand{\R}{\mathbb R}
\newcommand{\Q}{\mathbb Q}
\newcommand{\Z}{\mathbb Z}
\newcommand{\N}{\mathbb N}
\newcommand{\supp}{\operatorname{supp}}
\newcommand{\rad}{\operatorname{rad}}
\newcommand{\ord}{\operatorname{ord}}
\newcommand{\rank}{\operatorname{rank}}
\newcommand{\dist}{\operatorname{dist}}
\newcommand{\Span}{\operatorname{span}}
\newcommand{\vphi}{\varphi}
\newcommand{\vtheta}{\vartheta}
\newcommand{\val}{v}
\newcommand{\Dlog}{\mathscr D}
\newcommand{\newtag}{\textcolor{newgreen}{\textsf{NEW}}}
\newcommand{\inheritedtag}{\textcolor{bluegray}{\textsf{INHERITED}}}
\newcommand{\classtag}{\textcolor{navy}{\textsf{CLASSICAL}}}
\newcommand{\opentag}{\textcolor{openorange}{\textsf{OPEN}}}

\lstset{
  basicstyle=\ttfamily\small,
  columns=fullflexible,
  breaklines=true,
  frame=single,
  framerule=0.4pt,
  rulecolor=\color{bluegray},
  backgroundcolor=\color{softgray},
  showstringspaces=false
}

\title{\textbf{The Alaoglu--Erd\H{o}s Integer-Exponent Conjecture}\\[0.4em]
\Large Dyadic Cyclotomic Defects, Fresh Prime Supports,\\
and a Critical Matrix-Logarithm Target}
\author{Research continuation prepared by GPT-5.6 Pro}
\date{22 August 2026}

\begin{document}
\maketitle

\begin{center}
\fbox{\parbox{0.91\textwidth}{\centering
\textbf{Mathematical status.} This report proves several new exact lemmas and a sharp
no-go theorem, but it does \emph{not} prove the Alaoglu--Erd\H{o}s conjecture.  The final
support-sensitive determinant estimate identified below remains open.}}
\end{center}

\begin{abstract}
Assume that $x>0$ and that
\[
  M=2^x\in\Z,\qquad A=3^x\in\Z.
\]
Writing $\vtheta=\log 3/\log 2$, so that $A=M^{\vtheta}$, we introduce moving
logarithmic determinants built from $M^n\pm1$ and $A^n\pm1$.  Their signs are exact,
their magnitudes are exponentially small, and they satisfy positive telescoping
partitions indexed by every prime-power cyclotomic ladder.  In the dyadic ladder the
factors $M^{2^j}+1$ have pairwise disjoint odd prime supports; every odd divisor has
order exactly $2^{j+1}$, giving an infinite family of fresh-support, almost rank-one
$2\times2$ matrices of logarithms.  The defect has the exact interpretation of a
R\'enyi entropy and converges at the critical normalized slope
$\log2/\log3$.

We then prove a sharp contact ceiling: no fixed monic polynomial transfer can create a
smaller determinant, relative to logarithmic height, than the dyadic binomial family;
within the cyclotomic family, powers of two are uniquely optimal.  Finally, in the
rank-three external-support branch inherited from the unified report, we derive a new
integer certificate
\[
 D=a_2r^2-(m_2-a_3)rs-m_3s^2\ne0
\]
and an exact factorization connecting its sign to the rational approximation
$r\log3-s\log2$.

These results isolate a precise possible closing theorem: a quantitative lower bound for
matrices of logarithms that improves at the critical slope when the primitive output pair
has external prime support.  Qualitative Baker and Six Exponentials theorems do not supply
that improvement, and ordinary integer solutions exhibit the same analytic and
cyclotomic identities.  The report therefore separates rigorous progress from the still
missing arithmetic input and gives a concrete Lean formalization plan.
\end{abstract}

\tableofcontents
\newpage

\section{Executive assessment}

\begin{aeexecutivebox}
The new route is not another reformulation of the original determinant.  It manufactures
an infinite sequence of \emph{nonzero}, sign-controlled determinants whose columns acquire
fresh prime support at every dyadic level while their distance from rank one decays doubly
exponentially in the level.  The route reaches a sharp boundary: the same critical decay
occurs for every known integer solution $x\in\Z_{>0}$.  Consequently, the final theorem
must distinguish structural pairs $(2^m,3^m)$ from simultaneously power-primitive
pairs with primes outside $\{2,3\}$.
\end{aeexecutivebox}

The principal rigorous contributions of this continuation are as follows.

\begin{enumerate}
\item \textbf{Signed shifted determinants.}  The two families
\[
 (\log3)\log(M^n\pm1)-(\log2)\log(A^n\pm1)
\]
have opposite fixed signs, explicit two-sided bounds, and exact asymptotics.

\item \textbf{Positive prime-power partitions.}  For every prime $p$, the minus-defect
splits into a convergent sum of positive logarithmic determinants of
$\Phi_{p^j}(M^n)$ and $\Phi_{p^j}(A^n)$.  The partition exponentiates to an exact infinite
product independent of $p$.

\item \textbf{Entropy identity.}  Every geometric-quotient defect is exactly
$(\vtheta-1)\log2$ times a R\'enyi entropy of a finite geometric distribution.

\item \textbf{Dyadic support theorem.}  The odd parts of $T^{2^j}+1$ are pairwise
coprime.  Every odd prime divisor is private to its level and is congruent to $1$ modulo
$2^{j+1}$.  Hence the associated logarithms are rationally independent.

\item \textbf{Fresh-support near-rank-one matrices.}  The matrix with columns
$(\log2,\log3)^T$ and $(\log(M^{2^j}+1),\log(A^{2^j}+1))^T$ has a nonzero determinant
with exact sign, fresh arithmetic support, and critical decay rate $\log2/\log3$ relative
to height.

\item \textbf{Polynomial contact ceiling.}  A fixed monic polynomial of degree $d$ whose
first lower term is $cT^{d-r}$ has normalized contact slope $r/(d\vtheta)$.  This never
exceeds $1/\vtheta$; equality forces a binomial $T^d+c$.  For cyclotomic polynomials the
slope is $1/(\vtheta\vphi(\rad m))$, so powers of two are uniquely optimal.

\item \textbf{Rank-three integer certificate.}  When all external prime exponents lie on
one ray, the quadratic rank condition gives an explicit nonzero integer $D$, exact sign
information, and a lower bound for $|r\vtheta-s|$.
\end{enumerate}

\begin{aewarningbox}
None of the preceding statements rules out a noninteger $x$.  In particular, the private
prime theorem does not control overlap between the $M$-ladder and the $A$-ladder, and Six
Exponentials gives transcendence without a quantitative lower bound of the required
critical strength.
\end{aewarningbox}

\section{Source audit, current claims, and nonduplication}

\subsection{The attached unified report}

The source basis was the attached
\nolinkurl{alaoglu-erdos-unified-report.tex.gz}, a gzip-compressed LaTeX file.  After
decompression it contains $20{,}483$ lines ($974{,}791$ bytes).

The unified report already develops, among much else, the four-exponentials reduction,
exact conditional semigroup structure, primitive-generator reductions, valuation-lattice
geometry, determinant and interpolation hierarchies, $p$-adic and matrix-coefficient
approaches, entropy and automata viewpoints, and a quadratic-rank classification.  This
continuation therefore does \emph{not} repeat those surveys or their proofs.  It imports
only the minimal setup in \cref{sec:minimal-setup} and the rank-three external-ray
conclusion in \cref{sec:rank-three}.

A targeted text audit found no treatment in the source of the exact shifted families
$M^n\pm1$, their positive cyclotomic defect partitions, the dyadic pairwise-gcd/private-
prime theorem used here, or the fixed-polynomial contact ceiling.  The present report is
organized around those new items.

\subsection{Audit of the recent public claims}

The user's motivating claims were checked against public sources available on
22 August 2026.  The careful versions are recorded in \cref{tab:current-claims}.

\begin{table}[htbp]
\centering
\footnotesize
\begin{tabularx}{\textwidth}{@{}
>{\raggedright\arraybackslash}p{0.18\textwidth}
>{\raggedright\arraybackslash}X
>{\raggedright\arraybackslash}p{0.25\textwidth}@{}}
\toprule
Claim & Publicly supported statement & Necessary qualification \\
\midrule
Jacobian conjecture & A public Lean repository verifies an explicit polynomial map in
three variables with constant Jacobian determinant $-2$ and a three-point fiber.  This
refutes the usual conjecture in dimension $3$, hence in every dimension at least $3$. &
The two-dimensional conjecture remains open.  ``Resolved'' without this dimensional
qualification is misleading. \\
\addlinespace
Hadamard orders below $2000$ & OEIS A007299 records 2026 constructions for all twelve
previously unknown multiples of four at most $2000$: 668, 716, 892, 1132, 1244,
1388, 1436, 1676, 1772, 1916, 1948, and 1964. & The general Hadamard conjecture remains open. \\
\addlinespace
Elliptic curve of rank at least $30$ & The ICARM leaderboard lists curve \#273, submitted
20 August 2026, with a certified lower bound $\ge30$ and thirty independent rational
points. & The unconditional statement is a lower bound.  The page notes exact rank $30$
only under GRH and BSD. \\
\addlinespace
Alaoglu--Erd\H{o}s breakthrough & Targeted searches of current news, arXiv, and public
code repositories did not locate a public preprint, proof, or Lean repository matching
the private claim. & Absence from public search is not evidence that private work does not
exist.  There was nothing public to inspect or reproduce as of the access date. \\
\bottomrule
\end{tabularx}
\caption{Current-claim audit.  Source URLs and access dates appear in the bibliography.}
\label{tab:current-claims}
\end{table}

The Jacobian repository describes eleven Lean files with no \texttt{sorry} and verifies the
constant determinant and explicit collision.  ScienceDaily's republication of a
The Conversation article attributes the discovery to an AI-assisted workflow involving
Claude Fable 5, while explicitly saying that dimension two remains open.  OEIS attributes
the twelve Hadamard constructions to Alp\"oge, Philippe Voinov, Saul Reynolds-Haertle, and
Claude.  The ICARM commentary attributes the rank-$30$ construction to Claude with Levent
Alp\"oge and Ava Howell.  These attributions are reported as source statements rather than
independently reconstructed histories.

\subsection{Trust ledger}

\begin{aestatusbox}
\begin{itemize}
\item Statements proved from elementary real analysis, unique factorization, and finite
algebra are marked \newtag{} and supplied with proofs.
\item Baker's homogeneous linear-independence theorem and Six Exponentials are marked
\classtag{} and cited rather than reproved.
\item The rank-three external-ray reduction is marked \inheritedtag{} because it comes
from the attached unified report.
\item The quantitative support-sensitive matrix-logarithm estimate is marked \opentag{}.
No numerical experiment or heuristic is promoted to a proof.
\end{itemize}
\end{aestatusbox}

\section{Minimal inherited setup}\label{sec:minimal-setup}

The case $x<0$ is impossible because $0<2^x<1$ cannot be an integer.  The case $x=0$ is
already integral.  Hence throughout the main argument we assume
\begin{equation}\label{eq:basic}
 x>0,\qquad M=2^x\in\Z_{>1},\qquad A=3^x\in\Z_{>1}.
\end{equation}
Put
\begin{equation}\label{eq:theta-alpha}
 X=\log2,\qquad Y=\log3,\qquad
 \vtheta=\frac{Y}{X}=\log_2 3,
 \qquad \alpha=\frac1{\vtheta}=\frac{X}{Y}.
\end{equation}
Then
\begin{equation}\label{eq:power-curve}
 \log A=\vtheta\log M,
 \qquad A=M^{\vtheta}.
\end{equation}
Unique factorization shows that $\vtheta\notin\Q$.  In fact $\vtheta$ is transcendental:
if it were algebraic irrational, the Gelfond--Schneider theorem applied to
$2^{\vtheta}=3$ would give a contradiction.

For positive real numbers $B,C$, define the oriented logarithmic determinant
\begin{equation}\label{eq:Dlog}
 \Dlog(B,C)=Y\log B-X\log C
 =X\log\left(\frac{B^{\vtheta}}{C}\right).
\end{equation}
It is additive under componentwise multiplication:
\begin{equation}\label{eq:Dlog-add}
 \Dlog(B_1B_2,C_1C_2)=\Dlog(B_1,C_1)+\Dlog(B_2,C_2).
\end{equation}
The fundamental relation is simply
\begin{equation}\label{eq:Dlog-base-zero}
 \Dlog(M,A)=0,
 \qquad \Dlog(M^n,A^n)=0\quad(n\ge1).
\end{equation}

The attached unified report proves or records the following facts that will only be used
for context later.

\begin{itemize}
\item A hypothetical noninteger solution is transcendental and gives a primitive output
pair after passing to the least positive noninteger generator.
\item Its prime-exponent logarithmic tensor has quadratic rank $3$ or $4$; rank at most
$2$ is excluded by Baker's theorem.
\item In rank $3$, all exponent vectors at primes outside $\{2,3\}$ lie on one rational
ray.  This is the sole inherited structural input used in \cref{sec:rank-three}.
\end{itemize}

A pair $(M,A)$ will be called \emph{simultaneously power-primitive} when there is no
integer $d\ge2$ for which both $M$ and $A$ are perfect $d$th powers.  Equivalently, the
greatest common divisor of all nonzero prime valuations occurring in $M$ and $A$ is $1$.
The least nonintegral generator's output pair has this property.  We say that the pair has
\emph{external support} when a prime outside $\{2,3\}$ divides $MA$.

No statement in the next six sections requires the least-generator hypothesis: the exact
defect identities hold for every pair satisfying \eqref{eq:basic}.

\section{Signed shifted logarithmic defects}\label{sec:shifted}

The exact relation \eqref{eq:Dlog-base-zero} is unchanged by taking powers.  The first new
step is to perturb those powers by $\pm1$ before applying the logarithmic determinant.

\begin{definition}
For $n\ge1$, define
\begin{align}
 \Delta_n^+&=\Dlog(M^n+1,A^n+1),\label{eq:delta-plus}\\
 \Delta_n^-&=\Dlog(M^n-1,A^n-1),\label{eq:delta-minus}\\
 E_n&=-\Delta_n^- .\label{eq:E-def}
\end{align}
\end{definition}

\begin{theorem}[Sign, monotonicity, asymptotics, and dyadic conservation]
\label{thm:shifted-master}
For every $n\ge1$ one has
\begin{equation}\label{eq:signs}
 \Delta_n^-<0<\Delta_n^+,
 \qquad E_n>0.
\end{equation}
The sequence $(E_n)_{n\ge1}$ is strictly decreasing and tends to $0$.  Moreover,
\begin{align}
 E_n&=Y M^{-n}-X A^{-n}+O(M^{-2n}),\label{eq:E-asymp}\\
 \Delta_n^+&=Y M^{-n}-X A^{-n}+O(M^{-2n}),\label{eq:plus-asymp}
\end{align}
so that
\begin{equation}\label{eq:nth-root}
 \lim_{n\to\infty}E_n^{1/n}
 =\lim_{n\to\infty}(\Delta_n^+)^{1/n}=M^{-1}.
\end{equation}
Finally,
\begin{equation}\label{eq:dyadic-one-step}
 \Delta_n^++\Delta_n^-=\Delta_{2n}^-,
 \qquad
 \Delta_n^+=E_n-E_{2n},
\end{equation}
and hence the positive dyadic partition
\begin{equation}\label{eq:dyadic-sum}
 E_n=\sum_{j=0}^{\infty}\Delta_{2^j n}^+.
\end{equation}
\end{theorem}

\begin{proof}
Set $u=M^{-n}\in(0,1)$.  By \eqref{eq:power-curve}, $A^{-n}=u^{\vtheta}$.  Cancellation of
the leading power terms gives the exact formulas
\begin{align}
 \frac{\Delta_n^+}{X}
 &=\vtheta\log(1+u)-\log(1+u^{\vtheta}),\label{eq:plus-u}\\
 \frac{\Delta_n^-}{X}
 &=\vtheta\log(1-u)-\log(1-u^{\vtheta}).\label{eq:minus-u}
\end{align}
Since $\vtheta>1$ and $u>0$,
\[
 (1+u)^{\vtheta}>1+u^{\vtheta},
\]
which proves $\Delta_n^+>0$.  Also
\[
 u^{\vtheta}+(1-u)^{\vtheta}<u+(1-u)=1,
\]
so $(1-u)^{\vtheta}<1-u^{\vtheta}$ and $\Delta_n^-<0$.

For monotonicity, write
\[
 \frac{E_n}{X}=e(u):=-\vtheta\log(1-u)+\log(1-u^{\vtheta}).
\]
A direct derivative calculation yields
\begin{equation}\label{eq:e-derivative}
 e'(u)=
 \frac{\vtheta(1-u^{\vtheta-1})}{(1-u)(1-u^{\vtheta})}>0.
\end{equation}
As $n$ increases, $u=M^{-n}$ strictly decreases, and therefore $E_n$ strictly decreases.
The limit is $0$ by \eqref{eq:minus-u}.

Taylor expansion at $u=0$ gives
\[
 e(u)=\vtheta u-u^{\vtheta}+O(u^2),
 \qquad
 \frac{\Delta_n^+}{X}=\vtheta u-u^{\vtheta}+O(u^2),
\]
because $\vtheta>1$ and $2\vtheta>2$.  Multiplying by $X$, and using
$X\vtheta=Y$, proves \eqref{eq:E-asymp}--\eqref{eq:plus-asymp}; the root limits follow.

Finally, add the defining logarithms and use
$(M^n-1)(M^n+1)=M^{2n}-1$, and similarly for $A$:
\[
 \Delta_n^-+\Delta_n^+=\Delta_{2n}^-.
\]
This is \eqref{eq:dyadic-one-step}.  Iteration gives
\[
 E_n=\sum_{j=0}^{J}\Delta_{2^jn}^++E_{2^{J+1}n},
\]
and the remainder tends to $0$.
\end{proof}

The preceding asymptotics can be replaced by explicit inequalities.

\begin{proposition}[Sharp elementary two-sided bound]\label{prop:two-sided-plus}
For every $n\ge1$,
\begin{equation}\label{eq:plus-bound}
 Y M^{-n}
 \frac{1-(M/A)^n}{(1+M^{-n})(1+A^{-n})}
 \le \Delta_n^+
 <Y M^{-n}.
\end{equation}
Consequently,
\begin{equation}\label{eq:normalized-plus}
 \frac{\Delta_n^+M^n}{Y}\longrightarrow1.
\end{equation}
\end{proposition}

\begin{proof}
Let
\[
 f(s)=\vtheta\log(1+s)-\log(1+s^{\vtheta}).
\]
Then
\begin{equation}\label{eq:fprime}
 f'(s)=\frac{\vtheta(1-s^{\vtheta-1})}{(1+s)(1+s^{\vtheta})}.
\end{equation}
The numerator decreases and the denominator increases on $(0,1)$, so $f'$ decreases.
Thus, for $0<u<1$,
\[
 u f'(u)\le \int_0^u f'(s)\,ds=f(u)<\vtheta u.
\]
Substitute $u=M^{-n}$, note that $u^{\vtheta}=A^{-n}$ and
$u^{\vtheta-1}=(M/A)^n$, and multiply by $X$.
\end{proof}

\begin{remark}[A precise analytic picture]
The plus and minus defects have the same first-order magnitude, but the dyadic identity
shows that the plus defect is the amount by which the positive ``minus energy'' drops
when the scale doubles.  Thus $E_n$ is not merely asymptotic to a geometric sequence; it
is exactly the total mass of a positive dyadic cascade.
\end{remark}

\section{Geometric quotients, R\'enyi entropy, and cyclotomic partitions}
\label{sec:cyclotomic}

The dyadic identity is one instance of a prime-power factorization.  The correct finite
object is the geometric quotient
\begin{equation}\label{eq:Qkn}
 Q_{k,n}(T)=\frac{T^{kn}-1}{T^n-1}
 =\sum_{\ell=0}^{k-1}T^{\ell n}
 \qquad(k\ge2,n\ge1).
\end{equation}

\begin{definition}
Define the geometric-quotient defect
\begin{equation}\label{eq:Gkn}
 \mathcal G_{k,n}=\Dlog(Q_{k,n}(M),Q_{k,n}(A)).
\end{equation}
\end{definition}

\begin{theorem}[Exact positive quotient law]\label{thm:geometric-quotient}
For all $k\ge2$ and $n\ge1$,
\begin{equation}\label{eq:G-telescope}
 \mathcal G_{k,n}=\Delta_{kn}^--\Delta_n^-=E_n-E_{kn}>0.
\end{equation}
For fixed $n$, the sequence $\mathcal G_{k,n}$ increases strictly with $k$ and
\begin{equation}\label{eq:G-limit}
 \lim_{k\to\infty}\mathcal G_{k,n}=E_n.
\end{equation}
\end{theorem}

\begin{proof}
Apply the additivity \eqref{eq:Dlog-add} to the quotient in \eqref{eq:Qkn}.  This gives the
first equality in \eqref{eq:G-telescope}; the remainder follows from the strict decrease of
$E_n$.
\end{proof}

There is also a direct convexity proof with a useful information-theoretic interpretation.

\begin{theorem}[Exact R\'enyi-entropy identity]\label{thm:renyi}
Let
\begin{equation}\label{eq:geometric-prob}
 p_\ell=\frac{M^{\ell n}}{Q_{k,n}(M)}
 \qquad(0\le\ell<k).
\end{equation}
Then $(p_0,\dots,p_{k-1})$ is a probability vector and
\begin{equation}\label{eq:renyi-identity}
 \frac{\mathcal G_{k,n}}{X}
 =-\log\sum_{\ell=0}^{k-1}p_\ell^{\vtheta}
 =(\vtheta-1)H_{\vtheta}(p_0,\dots,p_{k-1}),
\end{equation}
where
\[
 H_{\vtheta}(p)=\frac{1}{1-\vtheta}
 \log\sum_\ell p_\ell^{\vtheta}
\]
is the R\'enyi entropy of order $\vtheta$.  In particular positivity is strict because the
probability vector has at least two positive coordinates.
\end{theorem}

\begin{proof}
Let $S=\sum_{\ell=0}^{k-1}M^{\ell n}=Q_{k,n}(M)$.  Since
$A^{\ell n}=(M^{\ell n})^{\vtheta}$,
\begin{align*}
 \frac{\mathcal G_{k,n}}X
 &=\vtheta\log S-
 \log\sum_{\ell=0}^{k-1}(M^{\ell n})^{\vtheta}\\
 &=-\log\sum_{\ell=0}^{k-1}
 \left(\frac{M^{\ell n}}S\right)^{\vtheta}.
\end{align*}
This is exactly \eqref{eq:renyi-identity}.
\end{proof}

\begin{remark}
The entropy identity does not add randomness to the problem.  It packages a deterministic
strict $\ell^1$--$\ell^{\vtheta}$ inequality.  Its value is that every prime-power
cyclotomic increment below becomes a positive entropy packet with an exact total budget.
\end{remark}

For a prime $p$ and $j\ge1$, define
\begin{equation}\label{eq:Cpjn}
 \mathcal C_{p^j,n}
 =\Dlog\bigl(\Phi_{p^j}(M^n),\Phi_{p^j}(A^n)\bigr).
\end{equation}
The prime-power identity
\begin{equation}\label{eq:prime-power-phi}
 \Phi_{p^j}(T^n)
 =\frac{T^{p^jn}-1}{T^{p^{j-1}n}-1}
 =\sum_{\ell=0}^{p-1}T^{\ell p^{j-1}n}
\end{equation}
places these factors inside \cref{thm:geometric-quotient}.

\begin{theorem}[Positive cyclotomic ladder]\label{thm:cyclotomic-ladder}
For every prime $p$, every $j\ge1$, and every $n\ge1$,
\begin{equation}\label{eq:C-difference}
 \mathcal C_{p^j,n}=E_{p^{j-1}n}-E_{p^jn}>0.
\end{equation}
Moreover,
\begin{equation}\label{eq:C-asymp}
 \mathcal C_{p^j,n}\sim Y M^{-p^{j-1}n}
 \qquad(j\to\infty),
\end{equation}
and one has the exact positive partition
\begin{equation}\label{eq:p-partition}
 E_n=\sum_{j=1}^{\infty}\mathcal C_{p^j,n}.
\end{equation}
Equivalently,
\begin{equation}\label{eq:p-product}
 \prod_{j=1}^{\infty}
 \frac{\Phi_{p^j}(M^n)^{\vtheta}}{\Phi_{p^j}(A^n)}
 =\frac{1-A^{-n}}{(1-M^{-n})^{\vtheta}}.
\end{equation}
The sum and product are independent of the chosen prime $p$.
\end{theorem}

\begin{proof}
Equation \eqref{eq:C-difference} follows by applying \eqref{eq:Dlog-add} to
\eqref{eq:prime-power-phi}.  Positivity follows from the decrease of $E_n$, or directly
from
\[
 \left(\sum_{\ell=0}^{p-1}y_\ell\right)^{\vtheta}
 >\sum_{\ell=0}^{p-1}y_\ell^{\vtheta}
\]
for $y_\ell=M^{\ell p^{j-1}n}>0$.  The asymptotic follows from
\eqref{eq:E-asymp}, because $E_{p^jn}=o(E_{p^{j-1}n})$.  Summation telescopes and the tail
tends to zero, proving \eqref{eq:p-partition}.  Divide by $X$ and exponentiate to obtain
\eqref{eq:p-product}.
\end{proof}

\begin{corollary}[Cross-prime conservation]\label{cor:cross-prime}
For any primes $p,q$ and every $n\ge1$,
\begin{equation}\label{eq:cross-prime}
 \sum_{j\ge1}\mathcal C_{p^j,n}
 =\sum_{j\ge1}\mathcal C_{q^j,n}=E_n.
\end{equation}
\end{corollary}

\begin{aenovelbox}
The identities \eqref{eq:p-partition}--\eqref{eq:cross-prime} create infinitely many
positive decompositions of the same small determinant.  Each decomposition has a distinct
arithmetic support pattern.  The dyadic version is strongest because its factors have
pairwise disjoint odd prime supports and, as proved later, it also attains the optimal
height-normalized contact slope among all fixed cyclotomic transfers.
\end{aenovelbox}

\section{Dyadic private primes and multiplicative independence}
\label{sec:private-primes}

Set
\begin{equation}\label{eq:Wj}
 W_j(T)=T^{2^j}+1\qquad(T>1,\ j\ge1).
\end{equation}
The following elementary fact is stronger than merely invoking a primitive-divisor
theorem.

\begin{theorem}[Exact dyadic gcd law]\label{thm:dyadic-gcd}
For integers $T>1$ and distinct $i,j\ge1$,
\begin{equation}\label{eq:dyadic-gcd}
 \gcd(W_i(T),W_j(T))=
 \begin{cases}
  1,&T\text{ even},\\
  2,&T\text{ odd}.
 \end{cases}
\end{equation}
Thus the odd parts of the $W_j(T)$ are pairwise coprime.
\end{theorem}

\begin{proof}
Assume $i<j$ and let $d$ divide both factors.  Modulo $d$,
$T^{2^i}\equiv-1$.  Raising to the even power $2^{j-i}$ gives
$T^{2^j}\equiv1$, while the second divisibility gives
$T^{2^j}\equiv-1$.  Hence $d\mid2$.  Parity now gives the two cases.
\end{proof}

\begin{theorem}[Private order and congruence]\label{thm:private-order}
For every $T>1$ and $j\ge1$, the integer $W_j(T)$ has an odd prime divisor.  Every odd
prime $q\mid W_j(T)$ satisfies
\begin{equation}\label{eq:order-congruence}
 \ord_q(T)=2^{j+1},
 \qquad q\equiv1\pmod{2^{j+1}}.
\end{equation}
In particular, $q$ divides no $W_i(T)$ with $i\ne j$ and $q\ge2^{j+1}+1$.
\end{theorem}

\begin{proof}
If $T$ is even, $W_j(T)>1$ is odd.  If $T$ is odd, then
$T^{2^j}\equiv1\pmod8$, so $W_j(T)\equiv2\pmod8$ and $W_j(T)>2$; it therefore has an odd
prime divisor $q$.  For such a divisor,
$T^{2^j}\equiv-1\pmod q$.  The order divides $2^{j+1}$ but does not divide $2^j$, so it is
exactly $2^{j+1}$.  Fermat's theorem gives the congruence for $q$.  Privacy also follows
from \cref{thm:dyadic-gcd}, or directly from the order.
\end{proof}

\begin{corollary}[Multiplicative and logarithmic independence]
\label{cor:W-independence}
For every integer $T>1$, the family $\{W_j(T):j\ge1\}$ is multiplicatively independent.
Consequently, the real numbers $\{\log W_j(T):j\ge1\}$ are linearly independent over
$\Q$.
\end{corollary}

\begin{proof}
In a finite relation $\prod_jW_j(T)^{c_j}=1$ with $c_j\in\Z$, choose an index with
$c_j\ne0$ and an odd prime divisor private to that index.  Its valuation forces $c_j=0$,
a contradiction.  Clearing denominators converts a rational linear relation among the
logs into such a multiplicative relation.
\end{proof}

The congruence in \eqref{eq:order-congruence} gives a quantitative, though modest, support
budget.

\begin{corollary}[Support growth]\label{cor:support-growth}
For each $J\ge1$, the product $\prod_{j=1}^JW_j(T)$ has at least $J$ distinct odd prime
divisors $q_j$ with
\begin{equation}\label{eq:qj-lower}
 q_j\ge2^{j+1}+1.
\end{equation}
Hence
\begin{equation}\label{eq:prime-log-budget}
 \sum_{j=1}^J\log q_j
 \ge \sum_{j=1}^J\log(2^{j+1}+1)
 =\frac{\log2}{2}J^2+O(J).
\end{equation}
\end{corollary}

Apply these statements to
\begin{equation}\label{eq:UVj}
 U_j=M^{2^j}+1,
 \qquad V_j=A^{2^j}+1
 \qquad(j\ge1).
\end{equation}
Because every private odd prime is $1$ modulo $4$, it is neither $2$ nor $3$.

\begin{corollary}[Independence with the base logs]\label{cor:base-log-independence}
Each of the two infinite families
\begin{equation}\label{eq:independent-families}
 \{\log2,\log3,\log U_1,\log U_2,\ldots\},
 \qquad
 \{\log2,\log3,\log V_1,\log V_2,\ldots\}
\end{equation}
is linearly independent over $\Q$ in every finite subfamily.
\end{corollary}

\begin{proof}
A private odd prime at each dyadic level first kills the coefficient of the corresponding
$\log U_j$ or $\log V_j$.  The remaining relation between $\log2$ and $\log3$ is trivial
by unique factorization.
\end{proof}

\begin{remark}[What is not proved]
The private prime for $U_j$ is private only relative to the other $U_i$; it may divide a
factor $V_k$.  No cross-family disjointness is asserted.  Controlling the intersections
$\supp(U_j)\cap\supp(V_k)$ is one possible adelic continuation of this route.
\end{remark}

\section{Fresh-support matrices approaching rank one}
\label{sec:moving-matrices}

Let $r_j=2^j$.  Form the logarithm matrix
\begin{equation}\label{eq:Lj}
 L_j=
 \begin{pmatrix}
  X&\log U_j\\
  Y&\log V_j
 \end{pmatrix}
 \qquad(j\ge1),
\end{equation}
and orient its small determinant by
\begin{equation}\label{eq:Gammaj}
 \Gamma_j=Y\log U_j-X\log V_j=-\det L_j.
\end{equation}
Thus $\Gamma_j=\Delta_{r_j}^+$.

\begin{theorem}[Fresh-support near-rank-one theorem]\label{thm:fresh-rank-one}
For every $j\ge1$:
\begin{enumerate}[label=\textup{(\roman*)}]
\item $\det L_j<0$, so $L_j$ has rank $2$;
\item the second entries $\log U_j$ (respectively $\log V_j$) acquire a private odd prime
support at every level;
\item the determinant satisfies
\begin{equation}\label{eq:Gamma-bound}
 Y M^{-r_j}
 \frac{1-(M/A)^{r_j}}{(1+M^{-r_j})(1+A^{-r_j})}
 \le\Gamma_j<Y M^{-r_j};
\end{equation}
\item in particular,
\begin{equation}\label{eq:Gamma-asymp}
 \Gamma_j\sim Y M^{-r_j};
\end{equation}
\item if $H_j=\log V_j$, then
\begin{equation}\label{eq:critical-slope}
 \lim_{j\to\infty}\frac{-\log\Gamma_j}{H_j}
 =\frac{\log M}{\log A}
 =\frac{X}{Y}=\alpha;
\end{equation}
\item the relative power approximation is exact:
\begin{equation}\label{eq:relative-power}
 \frac{U_j^{\vtheta}}{V_j}
 =\exp\!\left(\frac{\Gamma_j}{X}\right)>1,
 \qquad
 \frac{U_j^{\vtheta}}{V_j}-1\sim\vtheta M^{-r_j}.
\end{equation}
\end{enumerate}
\end{theorem}

\begin{proof}
The determinant assertion and the bound are \cref{prop:two-sided-plus} with $n=r_j$.
The support statement follows from \cref{thm:private-order}.  The asymptotic follows from
the bound.  Since
\[
 -\log\Gamma_j=r_j\log M+O(1),
 \qquad
 H_j=r_j\log A+o(1),
\]
we obtain \eqref{eq:critical-slope}.  Finally \eqref{eq:Dlog} gives
\[
 \frac{\Gamma_j}{X}=\log(U_j^{\vtheta}/V_j),
\]
and the final asymptotic follows from $e^t-1\sim t$.
\end{proof}

The phrase ``near rank one'' has an exact Euclidean meaning.  Let
\[
 c=\begin{pmatrix}X\\Y\end{pmatrix},
 \qquad
 w_j=\begin{pmatrix}\log U_j\\\log V_j\end{pmatrix}.
\]

\begin{proposition}[Exact distance to the rank-one line]\label{prop:distance-line}
For every $j\ge1$,
\begin{equation}\label{eq:distance-line}
 \dist(w_j,\R c)=\frac{\Gamma_j}{\sqrt{X^2+Y^2}}.
\end{equation}
Moreover,
\begin{equation}\label{eq:angle-asymp}
 \sin\angle(c,w_j)
 =\frac{\Gamma_j}{\|c\|\,\|w_j\|}
 \asymp\frac{M^{-r_j}}{r_j}.
\end{equation}
\end{proposition}

\begin{proof}
The distance from a vector $w$ to the line spanned by a nonzero vector $c$ in
$\R^2$ is $|\det(c,w)|/\|c\|$.  Here the absolute determinant is $\Gamma_j$.
The angle formula is the same identity divided by $\|w_j\|$, and
$\|w_j\|\asymp r_j$.
\end{proof}

There is a useful finite-product form of the dyadic conservation law.  From
\[
 \prod_{j=0}^{J}(T^{2^jn}+1)=\frac{T^{2^{J+1}n}-1}{T^n-1}
\]
and \cref{thm:cyclotomic-ladder}, one obtains
\begin{equation}\label{eq:dyadic-product}
 \prod_{j=0}^{\infty}
 \frac{(M^{2^jn}+1)^{\vtheta}}{A^{2^jn}+1}
 =\frac{1-A^{-n}}{(1-M^{-n})^{\vtheta}}.
\end{equation}
Every factor on the left is greater than $1$, and its logarithm is a positive defect.
Thus the moving matrices form a genuinely arithmetic, positive approach to rank one,
not a sequence with cancellation between signs.

\subsection{An exact moving-determinant equivalence}

For arbitrary integers $M,A>1$, not initially assumed to satisfy
$Y\log M=X\log A$, define $U_j,V_j,L_j$ as above.  Then
\begin{equation}\label{eq:general-Gamma}
 \Gamma_j
 =r_j\bigl(Y\log M-X\log A\bigr)
 +Y\log(1+M^{-r_j})-X\log(1+A^{-r_j}).
\end{equation}

\begin{proposition}[Decay is equivalent to the original determinant relation]
\label{prop:decay-equivalent}
For fixed integers $M,A>1$, the following are equivalent:
\begin{enumerate}[label=\textup{(\alph*)}]
\item $Y\log M-X\log A=0$;
\item $\Gamma_j\to0$;
\item $|\Gamma_j|\le C\rho^{r_j}$ for some $C>0$ and $0<\rho<1$.
\end{enumerate}
When these hold, $\Gamma_j>0$ for every $j$ and the critical slope is
$\log M/\log A$.
\end{proposition}

\begin{proof}
Equation \eqref{eq:general-Gamma} shows that if the base determinant is nonzero, then
$\Gamma_j/r_j$ tends to that nonzero value.  This excludes (b) and (c).  If the base
determinant vanishes, the proof of \cref{thm:fresh-rank-one} applies verbatim.
\end{proof}

\begin{remark}[Why this reformulation is useful but not free]
The Alaoglu--Erd\H{o}s conjecture is equivalent to classifying all integer pairs $(M,A)$
for which the second column of these fresh-support matrices approaches the base-log line
(or, equivalently, $\det L_j\to0$): they should be precisely
$(M,A)=(2^m,3^m)$ with $m\in\Z_{>0}$.  The equivalence does not itself reduce the logical
difficulty.  It does, however, replace a single exact zero by an infinite arithmetic
sequence carrying new, level-by-level prime supports.
\end{remark}

\section{Six Exponentials and the critical quantitative target}
\label{sec:six-exp}

The dyadic support theorem makes the qualitative Six Exponentials theorem applicable in a
particularly clean way.

\begin{theorem}[Qualitative transcendence along every triple]\label{thm:six-exp-triple}
Choose any three distinct indices $i,j,k\ge1$.  At least one of
\begin{equation}\label{eq:three-powers}
 U_i^{\vtheta},\qquad U_j^{\vtheta},\qquad U_k^{\vtheta}
\end{equation}
is transcendental.
\end{theorem}

\begin{proof}
The numbers $1$ and $\vtheta$ are linearly independent over $\Q$, while
$\log U_i,\log U_j,\log U_k$ are linearly independent over $\Q$ by
\cref{cor:W-independence}.  Six Exponentials says that at least one of the six numbers
\[
 \exp(\log U_\ell)=U_\ell,
 \qquad
 \exp(\vtheta\log U_\ell)=U_\ell^{\vtheta}
 \quad(\ell\in\{i,j,k\})
\]
is transcendental.  The first three are integers, so at least one of the second three is
transcendental.
\end{proof}

This conclusion is compatible with the integer $V_j$.  Indeed
$U_j^{\vtheta}\ne V_j$ by \eqref{eq:relative-power}; the approximation is only relative.

\begin{proposition}[Relative closeness but absolute separation]\label{prop:relative-not-absolute}
One has
\begin{align}
 \frac{U_j^{\vtheta}-V_j}{V_j}
 &\sim\vtheta M^{-r_j},\label{eq:relative-error}\\
 U_j^{\vtheta}-V_j
 &\sim\vtheta M^{(\vtheta-1)r_j}.
 \label{eq:absolute-error}
\end{align}
Thus the relative error tends doubly exponentially to zero in $j$, while the absolute
error tends to infinity.
\end{proposition}

\begin{proof}
The first statement is \eqref{eq:relative-power}.  Since
$V_j\sim A^{r_j}=M^{\vtheta r_j}$, multiplication gives the second.
\end{proof}

\begin{audit}
The construction supplies no small absolute distance from $U_j^{\vtheta}$ to an integer:
the distinguished integer $V_j$ is already far away in absolute terms, and no claim is made
about a different nearest integer.  The certified small quantity is the logarithmic
determinant, equivalently the relative error.  A useful transcendence measure must therefore
engage that height-normalized quantity rather than an absolute approximation not furnished by
the construction.
\end{audit}

Put
\begin{equation}\label{eq:height-slope}
 \lambda_j=\frac{-\log|\det L_j|}{\log V_j}.
\end{equation}
Under a hypothetical solution, \cref{thm:fresh-rank-one} gives
$\lambda_j\to\alpha$.  This identifies the exact quantitative threshold.

\begin{conjecture}[Support-sensitive critical-slope rigidity]
\label{conj:critical-rigidity}
Let $M,A>1$ be integers.  Suppose that $(M,A)$ is simultaneously power-primitive and
has external support.  Then there exists $\varepsilon(M,A)>0$ such that, for infinitely
many $j$,
\begin{equation}\label{eq:desired-lower}
 |\det L_j|
 \ge \exp\bigl(-(\alpha-\varepsilon(M,A))\log V_j\bigr).
\end{equation}
Equivalently,
\begin{equation}\label{eq:desired-slope}
 \liminf_{j\to\infty}\lambda_j<\alpha.
\end{equation}
The only pairs allowed to attain the critical slope $\alpha$ while satisfying
$Y\log M=X\log A$ should be $(2^m,3^m)$.
\end{conjecture}

\begin{aeopenbox}
If \cref{conj:critical-rigidity} were proved, a noninteger solution would yield the exact
asymptotic $\lambda_j\to\alpha$ and simultaneously the strict improvement
\eqref{eq:desired-slope}, a contradiction.  The desired theorem must be
\emph{base-sensitive}: the known integer solutions $(M,A)=(2^m,3^m)$ attain equality at
the same slope and cannot satisfy a universal strict improvement.
\end{aeopenbox}

The available classical theorems miss this target in distinct ways.

\begin{longtable}{@{}p{0.24\textwidth}p{0.32\textwidth}p{0.36\textwidth}@{}}
\toprule
Input & What it supplies here & Why it does not close the argument \\
\midrule
\endhead
Baker linear independence & Excludes rational quadratic forms of rank at most two at
independent prime logarithms. & The determinant is bilinear in logarithms and has rank
three or four in the remaining branches. \\
\addlinespace
Six Exponentials & At least one $U_j^{\vtheta}$ in every triple is transcendental. & It is
qualitative and gives no lower bound for the relative error at slope $\alpha$. \\
\addlinespace
Classical linear forms in logarithms & Effective lower bounds when coefficients are
algebraic and the expression is linear in logarithms. & Here the ``coefficients'' $X,Y$
are themselves logarithms, producing a $2\times2$ logarithmic determinant. \\
\addlinespace
Qualitative matrix-of-logarithms rank theorems & Nonvanishing and rank restrictions under
suitable support hypotheses. & They do not quantify how close a moving matrix can be to
rank one as its entries and supports grow. \\
\addlinespace
Subspace-theorem gcd bounds & Subexponential bounds for gcds such as
$\gcd(M^n-1,A^n-1)$ when $M,A$ are multiplicatively independent. & The elementary private-
prime theorem guarantees only $q_j\ge2^{j+1}+1$; subexponential gcd control still permits
shared prime mass of size $\exp(o(2^j))$ and does not force cross-family disjointness. \\
\bottomrule
\end{longtable}

\section{A sharp fixed-polynomial contact ceiling}
\label{sec:polynomial-ceiling}

The dyadic binomial may look arbitrary.  The next theorem shows that it is extremal among
all fixed monic polynomial transfers.

\begin{theorem}[Fixed-polynomial contact law]\label{thm:polynomial-contact}
Let $P\in\Z[T]$ be monic of degree $d\ge1$.  Suppose its first nonleading term is
\begin{equation}\label{eq:P-first-term}
 P(T)=T^d+cT^{d-r}+O(T^{d-r-1}),
 \qquad c\ne0,
 \qquad 1\le r\le d.
\end{equation}
For all sufficiently large $n$, $P(M^n)$ and $P(A^n)$ are positive.  Define
\begin{equation}\label{eq:Delta-P}
 \Delta_{P,n}=\Dlog(P(M^n),P(A^n)).
\end{equation}
Then
\begin{equation}\label{eq:Delta-P-asymp}
 \Delta_{P,n}
 =cY M^{-rn}-cX A^{-rn}+O(M^{-(r+1)n}),
\end{equation}
so, in particular,
\begin{equation}\label{eq:Delta-P-main}
 \Delta_{P,n}\sim cY M^{-rn}.
\end{equation}
Relative to the output height,
\begin{equation}\label{eq:contact-slope}
 \lim_{n\to\infty}
 \frac{-\log|\Delta_{P,n}|}{\log P(A^n)}
 =\frac{r}{d\vtheta}
 \le\frac1{\vtheta}=\alpha.
\end{equation}
Equality holds if and only if
\begin{equation}\label{eq:binomial-only}
 P(T)=T^d+c.
\end{equation}
\end{theorem}

\begin{proof}
Factor the leading term:
\[
 P(T)=T^d\left(1+cT^{-r}+O(T^{-r-1})\right).
\]
Because $2r\ge r+1$, the logarithm has the expansion
\begin{equation}\label{eq:logP}
 \log P(T)=d\log T+cT^{-r}+O(T^{-r-1}).
\end{equation}
Substitute $T=M^n$ and $T=A^n$.  The leading terms cancel by
$Y\log M=X\log A$, giving \eqref{eq:Delta-P-asymp}.  Since
$A^{-rn}=M^{-r\vtheta n}=o(M^{-rn})$, the main asymptotic follows.  Finally,
\[
 -\log|\Delta_{P,n}|=rn\log M+O(1),
 \qquad
 \log P(A^n)=dn\log A+o(n),
\]
which yields \eqref{eq:contact-slope}.  Since $1\le r\le d$, the slope is at most
$1/\vtheta$.  Equality forces $r=d$, meaning no term of degree $d-1,\ldots,1$ occurs;
thus $P(T)=T^d+c$.  The converse is immediate.
\end{proof}

\begin{aenovelbox}
The theorem is a rigorous no-go for a broad family of approximation ideas.  No fixed
polynomial with several lower terms can beat the critical dyadic slope.  To improve the
argument one must use changing arithmetic support, an adelic input, or a theorem that
classifies the equality case rather than merely optimizing the analytic expansion.
\end{aenovelbox}

Cyclotomic polynomials make the ceiling completely explicit.

\begin{lemma}[First lower cyclotomic term]\label{lem:cyclotomic-first}
Let $m\ge2$, $R=\rad(m)$, and $h=m/R$.  Then
\begin{equation}\label{eq:phi-rad}
 \Phi_m(T)=\Phi_R(T^h),
 \qquad
 \deg\Phi_m=h\vphi(R),
\end{equation}
and
\begin{equation}\label{eq:phi-first}
 \Phi_m(T)
 =T^{h\vphi(R)}-\mu(R)T^{h(\vphi(R)-1)}
 +O\!\left(T^{h(\vphi(R)-2)}\right),
\end{equation}
with the evident interpretation when $\vphi(R)=1$.
\end{lemma}

\begin{proof}
The identity \eqref{eq:phi-rad} is the standard reduction of a cyclotomic polynomial to
its squarefree kernel.  For squarefree $R$, the sum of the primitive $R$th roots of unity
is $\mu(R)$; therefore the coefficient below the leading coefficient of $\Phi_R$ is
$-\mu(R)$.  Substitute $T^h$.
\end{proof}

\begin{corollary}[Cyclotomic optimization]\label{cor:cyclotomic-optimal}
For fixed $m\ge2$,
\begin{equation}\label{eq:phi-defect-asymp}
 \Dlog(\Phi_m(M^n),\Phi_m(A^n))
 \sim-\mu(R)Y M^{-hn},
\end{equation}
where $R=\rad(m)$ and $h=m/R$.  Its height-normalized contact slope is
\begin{equation}\label{eq:phi-slope}
 \frac{1}{\vtheta\vphi(R)}.
\end{equation}
The maximum possible value is $1/\vtheta$, and among $m\ge2$ it occurs exactly when
$R=2$, i.e. when $m$ is a power of two.
\end{corollary}

\begin{proof}
Apply \cref{thm:polynomial-contact} with
$d=h\vphi(R)$, $r=h$, and $c=-\mu(R)$.  Since $R$ is squarefree and at least $2$,
$\vphi(R)\ge1$, with equality only for $R=2$.
\end{proof}

For $m=p^a$, the factor is a $p$-term geometric sum and its defect is positive by
\cref{thm:renyi}.  The dyadic case $p=2$ is therefore simultaneously:
\begin{itemize}
\item a positive entropy increment;
\item a factor with pairwise disjoint odd support across levels;
\item the unique cyclotomic family attaining the critical slope.
\end{itemize}
This triple optimality is the main reason to prioritize the dyadic ladder.

\section{The rank-three external-ray certificate}\label{sec:rank-three}

This section uses one structural statement from the attached unified report.  In its
rank-three branch, the exponent vectors at all primes outside $\{2,3\}$ lie on one rational
ray.  After taking the primitive integer direction of that ray, there is an integer
$C>1$ with $\gcd(C,6)=1$ and nonnegative integers
\begin{equation}\label{eq:external-ray-param}
 m_2,m_3,a_2,a_3,r,s
\end{equation}
with $r+s>0$ such that
\begin{equation}\label{eq:external-ray-MA}
 M=2^{m_2}3^{m_3}C^r,
 \qquad
 A=2^{a_2}3^{a_3}C^s.
\end{equation}
The case of a single external prime is obtained by taking $C=p$.

Put
\begin{equation}\label{eq:ZlogC}
 Z=\log C.
\end{equation}
Then $X,Y,Z$ are linearly independent over $\Q$.  Expanding
$Y\log M-X\log A=0$ gives the ternary quadratic relation
\begin{equation}\label{eq:Q-ray}
 Q(X,Y,Z)=
 -a_2X^2+(m_2-a_3)XY+m_3Y^2+(rY-sX)Z=0.
\end{equation}
Twice this quadratic form has symmetric matrix
\begin{equation}\label{eq:H-ray}
 H=
 \begin{pmatrix}
 -2a_2&m_2-a_3&-s\\
 m_2-a_3&2m_3&r\\
 -s&r&0
 \end{pmatrix}.
\end{equation}

\begin{theorem}[Nonzero integer rank certificate]\label{thm:D-certificate}
Define
\begin{equation}\label{eq:D-certificate}
 D=a_2r^2-(m_2-a_3)rs-m_3s^2\in\Z.
\end{equation}
Then
\begin{equation}\label{eq:detH}
 \det H=2D
\end{equation}
and
\begin{equation}\label{eq:D-nonzero}
 D\ne0.
\end{equation}
Furthermore, the exact identity
\begin{equation}\label{eq:D-factorization}
 DXY=(rY-sX)\bigl(a_2rX+m_3sY+rsZ\bigr)
\end{equation}
holds.
\end{theorem}

\begin{proof}
A direct determinant expansion gives \eqref{eq:detH}.  The quadratic form
\eqref{eq:Q-ray} is nonzero because $r+s>0$.  If $D=0$, then $H$ has rank at most $2$.
Over $\overline\Q$, a quadratic form of rank at most $2$ factors into two linear forms.
Its vanishing at the $\Q$-independent logarithms $X,Y,Z$ would then give a nontrivial
algebraic linear relation among those logarithms, contradicting Baker's homogeneous
linear-independence theorem.  Hence $D\ne0$.

For the factorization, a direct expansion gives the polynomial identity
\begin{align}
 &(rY-sX)(a_2rX+m_3sY+rsZ)-DXY\notag\\
 &\hspace{5em}=rs\,Q(X,Y,Z).
 \label{eq:factor-before-Q}
\end{align}
Equation \eqref{eq:Q-ray} now proves \eqref{eq:D-factorization}.  The same formula remains
valid when $r=0$ or $s=0$.
\end{proof}

\begin{corollary}[Edge branches and sign]\label{cor:D-branches}
Under the hypotheses above:
\begin{enumerate}[label=\textup{(\roman*)}]
\item if $r>0$ and $s=0$, then $a_2>0$, hence $2\mid A$;
\item if $r=0$ and $s>0$, then $m_3>0$, hence $3\mid M$;
\item if $r,s>0$, then
\begin{equation}\label{eq:sign-D}
 \operatorname{sgn}D=\operatorname{sgn}(r\log3-s\log2)
 =\operatorname{sgn}(r\vtheta-s).
\end{equation}
\end{enumerate}
\end{corollary}

\begin{proof}
If $s=0$, then $D=a_2r^2\ne0$, so $a_2>0$.  If $r=0$, then
$D=-m_3s^2\ne0$, so $m_3>0$.  If $r,s>0$, the second factor in
\eqref{eq:D-factorization} is positive, as are $X,Y$.
\end{proof}

Because $D$ is a nonzero integer, \eqref{eq:D-factorization} also yields a concrete
approximation lower bound.

\begin{corollary}[External-ray approximation bound]\label{cor:ray-approx}
If $r,s>0$, then
\begin{equation}\label{eq:ray-approx-raw}
 |rY-sX|
 =\frac{|D|XY}{a_2rX+m_3sY+rsZ}
 \ge \frac{XY}{r\log A+s\log M}.
\end{equation}
Equivalently,
\begin{equation}\label{eq:ray-theta-bound}
 |r\vtheta-s|
 \ge \frac{Y}{r\log A+s\log M}
 =\frac{\vtheta}{x(r\vtheta+s)}.
\end{equation}
\end{corollary}

\begin{proof}
The first equality is \eqref{eq:D-factorization}.  Since $|D|\ge1$ and
\begin{align*}
 r\log A+s\log M
 &=a_2rX+a_3rY+rsZ+m_2sX+m_3sY+rsZ\\
 &\ge a_2rX+m_3sY+rsZ,
\end{align*}
we obtain the inequality.  Divide by $X$ and use
$\log M=xX$, $\log A=xY$, and $Y=\vtheta X$.
\end{proof}

\begin{remark}[Interpretation]
The certificate $D$ is the exact integer obstruction preventing the rank-three quadratic
form from degenerating to rank two.  Its sign chooses the side from which $s/r$
approximates $\vtheta$, and its nonzero integrality gives a quantitative bound without any
unspecified constant.  The bound still contains the unknown $x$ and is not strong enough
to contradict continued-fraction approximation; it is a branch invariant, not a proof.
\end{remark}

\begin{aenovelbox}
The rank-three reduction in the unified report leaves a split ternary conic.  The new
factorization \eqref{eq:D-factorization} turns its nondegeneracy into a scalar integer
certificate and separates the edge cases $r=0$ and $s=0$.  This is suitable for a short
Lean formalization and for branch-specific searches.
\end{aenovelbox}

\section{Stress tests and exact failure points}\label{sec:stress-tests}

A useful progress report should try to break its own strongest-looking ideas.  The
following tests prevent the cyclotomic construction from being mistaken for a proof.

\subsection{Integer solutions are exact controls}

For every integer $m\ge1$, the pair
\begin{equation}\label{eq:integer-control}
 M=2^m,
 \qquad A=3^m
\end{equation}
satisfies every identity in \cref{sec:shifted,sec:cyclotomic,sec:private-primes,sec:moving-matrices}.
In particular:
\begin{itemize}
\item all defects have the same signs;
\item every prime-power partition is positive;
\item the dyadic factors have fresh private primes;
\item the moving logarithm matrices have critical slope $\alpha$;
\item Six Exponentials still forces transcendence among triples of the numbers
$(2^{m2^j}+1)^{\vtheta}$.
\end{itemize}
Thus none of those properties alone distinguishes integral from nonintegral $x$.

\begin{audit}[Control-family obstruction]
Any proposed determinant lower bound that is uniform over all $M,A$ and improves the
critical exponent below $\alpha$ is false, because it would also exclude the genuine
integer pairs \eqref{eq:integer-control}.  The arithmetic hypothesis must explicitly use
simultaneously power-primitive external support or must classify equality cases.
\end{audit}

\subsection{Fresh primes do not yet couple the two ladders}

The dyadic gcd law controls
\[
 \gcd(U_i,U_j),\qquad \gcd(V_i,V_j)
\]
for distinct levels, but not $\gcd(U_i,V_j)$.  If an odd prime $q\mid U_i$, then
$\ord_q(M)=2^{i+1}$.  If also $q\mid V_j$, then
$\ord_q(A)=2^{j+1}$.  The real relation
$\log A/\log M=\vtheta$ has no immediate meaning in the finite cyclic group
$(\Z/q\Z)^\times$.  Bridging that mismatch is precisely an adelic problem.

Known subspace-theorem estimates imply that for multiplicatively independent fixed
integers $M,A$, the gcd of $M^n-1$ and $A^n-1$ is subexponential in $n$.  The elementary
private-prime theorem gives only the lower bound $q_j\ge2^{j+1}+1$ and does not force a
prime of exponential size in $2^j$.  A subexponential gcd estimate therefore still permits
shared primes, or shared products of primes, of size $\exp(o(2^j))$.

\subsection{Qualitative transcendence is too coarse}

Six Exponentials proves that at least one $U_j^{\vtheta}$ in every triple is
transcendental, but the desired contradiction concerns the size of
\[
 \vtheta\log U_j-\log V_j.
\]
A transcendental number may be extremely close in relative terms to an integer.  The
classical theorem provides no exponent, no dependence on prime support, and no equality
classification at the slope $\alpha$.

\subsection{\texorpdfstring{$p$}{p}-adic order information is one-sided}

At a private prime $q_j\mid U_j$,
\[
 M^{2^j}\equiv-1\pmod{q_j},
 \qquad \ord_{q_j}(M)=2^{j+1}.
\]
This is exact and strong, but $q_j$ does not divide $M$, so the congruence does not directly
interact with the prime-exponent vector of $M$.  A successful $p$-adic step would need to
transport the real near-proportionality of the two log columns into a relation involving
$A$ modulo $q_j$ or a common resultant.  No such transport is currently proved.

\subsection{The polynomial ceiling blocks a common escape}

Replacing $T^{2^j}+1$ by a more elaborate fixed polynomial cannot improve the normalized
contact exponent.  \Cref{thm:polynomial-contact} proves that lower polynomial terms only
reduce the slope, and \cref{cor:cyclotomic-optimal} shows that all non-dyadic cyclotomic
families are worse by the factor $\vphi(\rad m)$.

\subsection{What remains genuinely plausible}

After the audits, three nonredundant directions remain.

\begin{enumerate}
\item \textbf{Cross-ladder adelic control.}  Use the private order conditions at divisors
of $U_j$ and $V_k$ together with resultants or support-sensitive matrix-logarithm bounds.

\item \textbf{Equality-case classification.}  Prove that critical determinant decay with
prime support outside $\{2,3\}$ is impossible, while explicitly allowing
$(2^m,3^m)$.

\item \textbf{Rank-three branch elimination.}  Combine the integer certificate $D$ with
additional local conditions on the external integer $C$.  The exact factorization is a
smaller target than the full rank-four problem.
\end{enumerate}

\section{Numerical calibration and reproducibility}\label{sec:numerics}

Numerics are included only to calibrate the exact asymptotics.  They are performed on the
known integer control $x=1$, where $M=2$ and $A=3$.  Let
\begin{equation}\label{eq:numerical-ratios}
 R_j=\frac{\Gamma_jM^{2^j}}{Y},
 \qquad
 \lambda_j=\frac{-\log\Gamma_j}{\log V_j}.
\end{equation}
The theorems predict $R_j\to1$ and $\lambda_j\to\alpha$.

\begin{table}[htbp]
\centering
\small
\begin{tabular}{@{}rrrrr@{}}
\toprule
$j$ & $2^j$ & $\Gamma_j$ & $R_j$ & $\lambda_j$ \\
\midrule
1 & 2   & $1.7211790324027\times10^{-1}$ & 0.626673868536 & 0.764173953182 \\
2 & 4   & $5.8097974438167\times10^{-2}$ & 0.846128885139 & 0.645746715285 \\
3 & 8   & $4.1774557447219\times10^{-3}$ & 0.973435926104 & 0.623281520363 \\
4 & 16  & $1.6747263074177\times10^{-5}$ & 0.999031818732 & 0.625634482844 \\
5 & 32  & $2.5579023220339\times10^{-10}$& 0.999998537502 & 0.628254606989 \\
6 & 64  & $5.9555891504460\times10^{-20}$& 0.999999999997 & 0.629592159480 \\
7 & 128 & $3.2285313476833\times10^{-39}$& 1.000000000000 & 0.630260956526 \\
\bottomrule
\end{tabular}
\caption{Dyadic determinant calibration for the integer control $x=1$.  The limiting
value is $\alpha=0.630929753571457\ldots$.}
\label{tab:numerical}
\end{table}

The nonmonotonic early behavior of $\lambda_j$ is harmless: the denominator
$\log(3^{2^j}+1)$ and the constant in $-\log\Gamma_j$ still influence small levels.
The normalized amplitude $R_j$ rapidly confirms the explicit bound
\eqref{eq:Gamma-bound}.

The following short script reproduces the table without subtracting large nearly equal
quantities; it uses \texttt{log1p} on the small tails.

\begin{lstlisting}[language=Python,caption={High-precision numerical calibration.}]
import mpmath as mp

mp.mp.dps = 100
x = mp.mpf(1)
M = mp.power(2, x)
A = mp.power(3, x)
X = mp.log(2)
Y = mp.log(3)
alpha = X / Y

for j in range(1, 8):
    r = 2**j
    gamma = Y*mp.log1p(M**(-r)) - X*mp.log1p(A**(-r))
    logV = r*mp.log(A) + mp.log1p(A**(-r))
    normalized = gamma*M**r/Y
    slope = -mp.log(gamma)/logV
    print(j, r, gamma, normalized, slope)

print("alpha =", alpha)
\end{lstlisting}

A symbolic check of the rank-three identity can be reduced to
\begin{lstlisting}[language=Python,caption={Symbolic verification of the certificate identity.}]
import sympy as sp
X,Y,Z,a2,m2,a3,m3,r,s = sp.symbols(
    "X Y Z a2 m2 a3 m3 r s")
Q = (-a2*X**2 + (m2-a3)*X*Y + m3*Y**2
     + (r*Y-s*X)*Z)
D = a2*r**2 - (m2-a3)*r*s - m3*s**2
R = (r*Y-s*X)*(a2*r*X + m3*s*Y + r*s*Z)
assert sp.factor(R - D*X*Y - r*s*Q) == 0
\end{lstlisting}

\section{Lean formalization blueprint}\label{sec:lean}

The new exact layer is deliberately elementary enough to formalize before attempting the
open quantitative theorem.  A modular plan is given in \cref{tab:lean-plan}.

\footnotesize
\begin{longtable}{@{}>{\raggedright\arraybackslash}p{0.20\textwidth}
>{\raggedright\arraybackslash}p{0.40\textwidth}
>{\raggedright\arraybackslash}p{0.32\textwidth}@{}}
\toprule
Proposed module & Principal theorem statements & Main Mathlib dependencies \\
\midrule
\endhead
\shortstack[l]{\ttfamily ShiftedLog\\\ttfamily Defect} & Exact $u$-formulas, signs of $\Delta_n^\pm$, derivative of
$E$, dyadic identity, explicit bound, asymptotics. & \texttt{Real.log}, derivatives,
real powers, monotonicity, finite/infinite geometric limits. \\
\addlinespace
\shortstack[l]{\ttfamily GeometricQuotient\\\ttfamily Defect} &
$\mathcal G_{k,n}=E_n-E_{kn}>0$ and the finite-sum convexity proof. & Finite sums,
\texttt{Real.rpow}, strict convexity or the power-sum inequality. \\
\addlinespace
\shortstack[l]{\ttfamily RenyiCyclotomic\\\ttfamily Defect} & Exact R\'enyi identity, prime-power cyclotomic
partition, telescoping sum and product. & Cyclotomic identities already in algebraic
number-theory libraries; analytic convergence. \\
\addlinespace
\shortstack[l]{\ttfamily DyadicPrivate\\\ttfamily Prime} & Exact gcd law, existence of an odd divisor, order
$2^{j+1}$, congruence $q\equiv1$, privacy. & \texttt{Nat.exists\_prime\_and\_dvd}, modular
powers, multiplicative order, elementary parity. \\
\addlinespace
\shortstack[l]{\ttfamily DyadicLog\\\ttfamily Independence} & Multiplicative independence and finite rational
independence of the logarithms. & Prime valuations, unique factorization, clearing
rational denominators. \\
\addlinespace
\shortstack[l]{\ttfamily MovingLog\\\ttfamily Matrix} & Determinant sign, distance-to-line formula, critical slope,
relative/absolute error split. & $2\times2$ matrices, Euclidean norm, elementary limits. \\
\addlinespace
\shortstack[l]{\ttfamily PolynomialContact\\\ttfamily Ceiling} & Log expansion for a fixed monic polynomial and the
slope $r/(d\vtheta)$. & Polynomial degree/coefficient API, eventual positivity,
Taylor estimates for \texttt{Real.log}. \\
\addlinespace
\shortstack[l]{\ttfamily ExternalRay\\\ttfamily Certificate} & $\det H=2D$, polynomial identity
$R-DXY=rsQ$, edge cases and sign consequences. & Pure ring normalization
(\texttt{ring}), integer nonzero arithmetic; Baker remains an explicit imported
mathematical hypothesis unless a suitable theorem is already formalized. \\
\bottomrule
\caption{A proposed Lean decomposition.}
\label{tab:lean-plan}
\end{longtable}
\normalsize

\subsection{Suggested theorem signatures}

The finite algebraic core can be separated from analytic assumptions.  Schematic Lean
interfaces are:

\begin{lstlisting}[language={},caption={Schematic Lean interfaces, not compiled code.}]
-- Pure determinant/cyclotomic algebra
 theorem shifted_add_eq_minus_double
   (M A n : Nat) (hM : 1 < M) (hA : 1 < A) :
   deltaPlus M A n + deltaMinus M A n = deltaMinus M A (2*n)

 theorem dyadic_gcd
   (T i j : Nat) (hT : 1 < T) (hij : i != j) (hi : 1 <= i) (hj : 1 <= j) :
   Nat.gcd (T^(2^i) + 1) (T^(2^j) + 1) = if Even T then 1 else 2

 theorem externalRay_ring_identity :
   (r*Y-s*X)*(a2*r*X+m3*s*Y+r*s*Z) - D*X*Y = r*s*Q

-- Analytic layer under the exact base-log relation
 theorem deltaPlus_pos
   (hrel : Real.log 3 * Real.log M = Real.log 2 * Real.log A) :
   0 < deltaPlus M A n
\end{lstlisting}

The exact gcd and ring identities should be formalized first: they are short, stable, and
independent of transcendence infrastructure.  The next priority is the sign theorem for
$\Delta_n^+$, because it packages the strict power-sum inequality that also proves all
prime-power cyclotomic increments positive.

\subsection{Formalization will not manufacture the missing theorem}

Lean can certify every exact identity and prevent silent sign or index errors.  It cannot
replace \cref{conj:critical-rigidity}.  A faithful final theorem should expose any new
matrix-logarithm or transcendence estimate as an explicit hypothesis until a proof is
available.  The report's recommended workflow is therefore:
\begin{enumerate}
\item formalize the elementary ladder and certificate modules;
\item use them to state the critical-slope theorem with no hidden assumptions;
\item search for the missing adelic input while the exact interface remains fixed;
\item only then connect the new theorem to the final contradiction.
\end{enumerate}

\section{Prioritized research program}\label{sec:program}

\subsection{Priority A: cross-ladder resultants}

For $i,j\ge1$, study
\begin{equation}\label{eq:cross-gcd}
 G_{i,j}=\gcd(M^{2^i}+1,A^{2^j}+1).
\end{equation}
An odd prime $q\mid G_{i,j}$ imposes simultaneous exact orders
\begin{equation}\label{eq:cross-orders}
 \ord_q(M)=2^{i+1},
 \qquad
 \ord_q(A)=2^{j+1}.
\end{equation}
Three concrete questions are now sharply posed.

\begin{enumerate}
\item Can a support-sensitive resultant bound force most private prime mass of $U_i$ to
remain outside every $V_j$ in a controlled range?
\item Does the simultaneously power-primitive external support of $(M,A)$ imply a lower bound for the product of
private prime parts not shared by the opposite ladder?
\item Can \eqref{eq:cross-orders} be combined with a $p$-adic logarithm estimate after
passing to suitable powers that are $1$ modulo $q$?
\end{enumerate}

A negative answer at one level is not enough; the determinant target needs an aggregate
bound over infinitely many levels.

\subsection{Priority B: equality cases in matrix-logarithm estimates}

Rather than seek a stronger generic lower bound, seek a theorem with a sharp exceptional
set.  A model statement is:

\begin{quote}
If an integer pair $(M,A)$ has the dyadic fresh-support matrices \eqref{eq:Lj} at critical
slope $\alpha$, then either $(M,A)=(2^m,3^m)$ or its external prime-exponent tensor satisfies
an explicit exceptional algebraic configuration.
\end{quote}

The current rank-three certificate suggests that the exceptional configuration may be
expressed through the integer $D$ and the external integer $C$.  Even a theorem reducing
all critical pairs to finitely many ray types would be substantial progress.

\subsection{Priority C: exploit all prime partitions simultaneously}

For every prime $p$, the same $E_n$ has the positive partition
\[
 E_n=\sum_{j\ge1}\mathcal C_{p^j,n}.
\]
The equality of the $p$- and $q$-partitions is analytically automatic, but their factor
supports differ.  A genuinely adelic invariant could compare the prime ideals appearing
in these two factorizations before taking logarithms.  The dyadic partition is optimal in
contact, while an odd-prime partition has more terms per factor and different order
constraints; using both may prevent the integer-control symmetry from remaining invisible.

\subsection{Priority D: finish the rank-three branch first}

In rank three, write $M,A$ as in \eqref{eq:external-ray-MA}.  The following finite data are
available:
\begin{equation}\label{eq:rank-three-data}
 (m_2,m_3,a_2,a_3,r,s,D,C),
 \qquad D\ne0,
 \qquad \operatorname{sgn}D=\operatorname{sgn}(r\vtheta-s).
\end{equation}
A branch search should combine:
\begin{itemize}
\item exact local restrictions on $C$ and on the parity of $M,A$;
\item continued-fraction information on $s/r$;
\item primitive-output conditions from the unified report;
\item private-prime orders in $M^{2^j}+1$ and $A^{2^j}+1$.
\end{itemize}
If rank three can be eliminated, the remaining rank-four problem will have a clearer
statement: genuinely two-dimensional external support is necessary.

\section{Conclusion}

The continuation has produced an exact new architecture rather than a claimed final proof.
Starting from a hypothetical integer point $(M,A)$ on $A=M^{\log_2 3}$, one obtains:

\begin{enumerate}
\item signed logarithmic defects at $M^n\pm1$ and $A^n\pm1$;
\item positive telescoping decompositions through every prime-power cyclotomic ladder;
\item an exact R\'enyi-entropy formula for each increment;
\item a dyadic family whose odd prime supports are pairwise disjoint;
\item fresh-support $2\times2$ logarithm matrices approaching rank one at the exact
critical slope $\alpha=\log2/\log3$;
\item a proof that fixed polynomial and cyclotomic transfers cannot improve that slope;
\item a nonzero integer certificate controlling the inherited rank-three external-ray
branch.
\end{enumerate}

The strongest conceptual advance is the separation of two phenomena that previously sat
inside a single zero determinant:
\begin{equation}\label{eq:final-separation}
 \text{analytic critical contact}
 \quad+\quad
 \text{level-by-level fresh arithmetic support}.
\end{equation}
The analytic contact alone is shared by all integer solutions.  The fresh support alone is
also shared by all integer solutions.  The missing theorem must couple the two while
recognizing $(2^m,3^m)$ as the exact equality family.

\begin{aeopenbox}
The most precise current target is \cref{conj:critical-rigidity}: prove a strict
height-normalized determinant improvement for simultaneously power-primitive pairs with external support, or
classify all equality cases.  This is narrower than a generic four-exponentials theorem,
because it concerns a very structured sequence of moving matrices with pairwise private
dyadic supports.  It is also honest: no currently cited theorem supplies the required
quantitative, support-sensitive conclusion.
\end{aeopenbox}

\appendix

\section{A compact proof ledger}\label{app:ledger}

\begin{longtable}{@{}p{0.35\textwidth}p{0.16\textwidth}p{0.39\textwidth}@{}}
\toprule
Statement & Status & Dependency \\
\midrule
\endhead
Signs and monotonicity of $\Delta_n^\pm,E_n$ & \newtag{} & Elementary real analysis \\
Positive dyadic and prime-power partitions & \newtag{} & Factorization and telescoping \\
R\'enyi identity & \newtag{} & Algebraic normalization of a finite sum \\
Dyadic exact gcd/private order & \newtag{} & Modular arithmetic and parity \\
Logarithmic independence of dyadic factors & \newtag{} & Unique factorization \\
Fresh-support near-rank-one theorem & \newtag{} & Previous elementary results \\
Transcendence in every triple & \classtag{} & Six Exponentials \\
Polynomial contact ceiling & \newtag{} & Fixed polynomial log expansion \\
Cyclotomic optimum & \newtag{} & Squarefree-kernel identity for $\Phi_m$ \\
External-ray reduction & \inheritedtag{} & Attached unified report \\
Nonzero certificate $D$ & \newtag{} + \classtag{} & Determinant algebra plus Baker \\
Support-sensitive critical-slope rigidity & \opentag{} & Missing quantitative theorem \\
\bottomrule
\caption{Proof and trust ledger.}
\end{longtable}

\section{Selected derivations in expanded form}\label{app:derivations}

\subsection{The dyadic infinite product}

For finite $J$,
\[
 \prod_{j=0}^{J}(1+t^{2^j})=\frac{1-t^{2^{J+1}}}{1-t}.
\]
Apply this first with $t=M^{-n}$ and then with $t=A^{-n}$.  Raising the first identity to
$\vtheta$, dividing by the second, and letting $J\to\infty$ gives
\[
 \prod_{j=0}^{\infty}
 \frac{(1+M^{-2^jn})^{\vtheta}}{1+A^{-2^jn}}
 =\frac{1-A^{-n}}{(1-M^{-n})^{\vtheta}}.
\]
The large powers cancel inside each ratio because
\[
 \frac{(M^{2^jn}+1)^{\vtheta}}{A^{2^jn}+1}
 =\frac{(1+M^{-2^jn})^{\vtheta}}{1+A^{-2^jn}}.
\]
This proves \eqref{eq:dyadic-product} without interchanging a conditionally convergent
series: every logarithmic term is positive and the finite products telescope exactly.

\subsection{The entropy packet at a prime-power level}

Let $h=p^{j-1}n$ and $s=M^{-h}$.  Reverse the $p$ terms in
$\Phi_{p^j}(M^n)$ and normalize:
\[
 \Phi_{p^j}(M^n)=M^{(p-1)h}(1+s+\cdots+s^{p-1}).
\]
The corresponding $A$ factor is
\[
 \Phi_{p^j}(A^n)=A^{(p-1)h}(1+s^{\vtheta}+\cdots+s^{(p-1)\vtheta}).
\]
The leading terms cancel in $\Dlog$, leaving
\[
 \frac{\mathcal C_{p^j,n}}X
 =\vtheta\log(1+s+\cdots+s^{p-1})
 -\log(1+s^{\vtheta}+\cdots+s^{(p-1)\vtheta}).
\]
If
\[
 \pi_\ell=\frac{s^\ell}{1+s+\cdots+s^{p-1}},
\]
then this is
\[
 -\log\sum_{\ell=0}^{p-1}\pi_\ell^{\vtheta}
 =(\vtheta-1)H_{\vtheta}(\pi).
\]
The first correction comes from $\ell=1$, so
$\mathcal C_{p^j,n}\sim Y s=Y M^{-p^{j-1}n}$.

\subsection{Why dyadic factors are optimally tangent}

For a fixed monic $P$, the first lower term determines the first nonzero perturbation of
$\log P(T)$ from $d\log T$.  A gap of $r$ degrees creates a tail $T^{-r}$, while the
height grows like $d\log T$.  Their ratio is therefore $r/d$.  The largest possible gap is
$r=d$, when the only lower term is constant.  A dyadic cyclotomic polynomial is exactly
such a binomial:
\[
 \Phi_{2^a}(T)=T^{2^{a-1}}+1.
\]
Every cyclotomic polynomial whose squarefree kernel contains an odd prime has at least
$\vphi(R)\ge2$ primitive-root directions, reducing the normalized gap by that factor.

\section{Bibliography and current-source ledger}

\begin{thebibliography}{99}

\bibitem{AlaogluErdos1944}
L.~Alaoglu and P.~Erd\H{o}s,
\newblock On highly composite and similar numbers,
\newblock \emph{Transactions of the American Mathematical Society} \textbf{56} (1944),
448--469.

\bibitem{UnifiedReport2026}
\emph{The Alaoglu--Erd\H{o}s Integer-Exponent Conjecture: Unified Report},
\newblock attached LaTeX source, audited 22 August 2026.

\bibitem{Baker1975}
A.~Baker,
\newblock \emph{Transcendental Number Theory},
\newblock Cambridge University Press, 1975.

\bibitem{Waldschmidt2000}
M.~Waldschmidt,
\newblock \emph{Diophantine Approximation on Linear Algebraic Groups},
\newblock Grundlehren der mathematischen Wissenschaften 326, Springer, 2000.

\bibitem{Lang1966}
S.~Lang,
\newblock \emph{Introduction to Transcendental Numbers},
\newblock Addison--Wesley, 1966.

\bibitem{LaurentMignotteNesterenko1995}
M.~Laurent, M.~Mignotte, and Yu.~Nesterenko,
\newblock Formes lin\'eaires en deux logarithmes et d\'eterminants d'interpolation,
\newblock \emph{Journal of Number Theory} \textbf{55} (1995), 285--321.

\bibitem{DasguptaKakde2026}
S.~Dasgupta and M.~Kakde,
\newblock Ranks of matrices of logarithms of algebraic numbers II: The Matrix
Coefficient Conjecture,
\newblock \emph{Advances in Mathematics} \textbf{487} (2026), Article 110753;
\newblock \url{https://doi.org/10.1016/j.aim.2025.110753}; arXiv:2408.08178.

\bibitem{BCZ2003}
Y.~Bugeaud, P.~Corvaja, and U.~Zannier,
\newblock An upper bound for the G.C.D. of $a^n-1$ and $b^n-1$,
\newblock \emph{Mathematische Zeitschrift} \textbf{243} (2003), 79--84.

\bibitem{HLP}
G.~Hardy, J.~E.~Littlewood, and G.~P\'olya,
\newblock \emph{Inequalities},
\newblock Cambridge University Press, second edition, 1952.

\bibitem{JacobianLean2026}
A.~Rader et al.,
\newblock \emph{Formal verification and structure theory of Alp\"oge's counterexample to
the Jacobian Conjecture}, Lean 4 repository,
\newblock \url{https://github.com/alerad/alpoge-lean}, accessed 22 August 2026.

\bibitem{ScienceDailyJacobian2026}
The Conversation / ScienceDaily,
\newblock Claude Fable 5 AI finds a tiny formula that topples an 87-year-old math
conjecture,
\newblock 5 August 2026,
\url{https://www.sciencedaily.com/releases/2026/08/260804034634.htm}, accessed
22 August 2026.

\bibitem{OEISHadamard2026}
OEIS Foundation,
\newblock A007299, comments updated 12 August 2026,
\newblock \url{https://oeis.org/A007299}, accessed 22 August 2026.

\bibitem{ICARMRank30}
ICARM Elliptic Curve Rank Leaderboard,
\newblock Curve \#273, submitted 20 August 2026,
\newblock \url{https://elliptic-rank.icarm.cloud/curve/273}, accessed 22 August 2026.

\end{thebibliography}

\end{document}
